\chapter{微分方程组}
前几章讨论的是一个未知函数所满足的微分方程. 在实际问题与理论研究中，
还经常出现多个未知函数相互耦合的情形，这便引出微分方程组.

本章可类比第三章，建立线性方程组的解空间理论，并讨论一些特殊情况的微
分方程组的求解方法.

\section{微分方程组的概念}\label{sec:4.1}

\subsection{基本概念}

\begin{deff}[(一阶微分方程组及其解与初值问题)]
含有 $n$ 个未知函数 $x_1(t),\ldots,x_n(t)$，并且每个方程关于未知函数的
导数都是一阶的方程组
\begin{equation}\label{equ:4.1}
\left\{
\begin{aligned}
&\frac{dx_1}{dt}=f_1(t,x_1,\ldots,x_n),\\
&\frac{dx_2}{dt}=f_2(t,x_1,\ldots,x_n),\\
&\qquad\vdots\\
&\frac{dx_n}{dt}=f_n(t,x_1,\ldots,x_n)
\end{aligned}
\right.
\end{equation}
称为\textbf{一阶微分方程组}.

方程组中出现的未知函数导数的最高阶数称为方程组的\textbf{阶数}. 若每个
$f_i(t,x_1,\ldots,x_n)$ 关于所有未知函数都是一次的，则称该方程组为
\textbf{线性微分方程组}；否则称为\textbf{非线性微分方程组}.

对于一阶线性方程组，若我们记
\[
\bm{x}=
\begin{pmatrix}
x_1\\ \vdots\\ x_n
\end{pmatrix},\quad
A(t)=
\begin{pmatrix}
a_{11}(t)&\cdots&a_{1n}(t)\\
\vdots&&\vdots\\
a_{n1}(t)&\cdots&a_{nn}(t)
\end{pmatrix},\quad
\bm{f}(t)=
\begin{pmatrix}
f_1(t)\\ \vdots\\ f_n(t)
\end{pmatrix},
\]
则 (\ref{equ:4.1}) 可写成矩阵形式
\begin{equation}\label{equ:4.2}
\frac{d\bm{x}}{dt}=A(t)\bm{x}+\bm{f}(t).
\end{equation}
当 $\bm{f}(t)\equiv\bm{0}$ 时称为\textbf{齐次线性方程组}，否则称为
\textbf{非齐次线性方程组}.

若函数组 $x_1(t),\ldots,x_n(t)$ 在区间 $I$ 上可微，并且代入
(\ref{equ:4.1}) 后在 $I$ 上使每个方程恒成立，则称其为方程组在
$I$ 上的一个\textbf{解}. 含有 $n$ 个相互独立任意常数
$C_1,\ldots,C_n$ 的解
\[
x_i=\varphi_i(t,C_1,\ldots,C_n),\qquad i=1,\ldots,n
\]
称为方程组的\textbf{通解}. 若由 $n$ 个关系式
\[
\Phi_i(t,x_1,\ldots,x_n;C_1,\ldots,C_n)=0,
\qquad i=1,\ldots,n
\]
隐式给出通解，则称其为\textbf{通积分}. 给定
\begin{equation}\label{equ:4.3}
x_1(t_0)=x_{10},\quad x_2(t_0)=x_{20},\quad\ldots,\quad x_n(t_0)=x_{n0},
\end{equation}
称为方程组的\textbf{初始条件}；方程组与初始条件合在一起称为
\textbf{初值问题}. 将初始条件代入通解或通积分确定常数后所得的解称为
\textbf{特解}.
\end{deff}

\subsection{高阶方程组化为一阶方程组}

高阶微分方程组可以通过增加状态变量化为一阶方程组.

\begin{prop}\label{prop:4.1.1}
设显式 $n$ 阶方程
\begin{equation}\label{equ:4.4}
y^{(n)}=f\left(t,y,y',\ldots,y^{(n-1)}\right).
\end{equation}
令
\[
x_1=y,\qquad x_2=y',\qquad\ldots,\qquad x_n=y^{(n-1)},
\]
则 (\ref{equ:4.4}) 等价于
\begin{equation}\label{equ:4.5}
\left\{
\begin{aligned}
x_1'&=x_2,\\
x_2'&=x_3,\\
&\vdots\\
x_{n-1}'&=x_n,\\
x_n'&=f(t,x_1,\ldots,x_n).
\end{aligned}
\right.
\end{equation}
\end{prop}

\begin{pf1}
前 $n-1$ 个等式由变量定义直接得到. 最后一个等式等价于
\[
x_n'=y^{(n)}=f(t,y,\ldots,y^{(n-1)}),
\]
也就是 $x_n'=f(t,x_1,\ldots,x_n)$. 反过来，由 (\ref{equ:4.5}) 的链式关系
可逐一得到 $x_i=y^{(i-1)}$，最后得到 (\ref{equ:4.4}).
\end{pf1}

上述结论可推广到高阶方程组. 设系统含有 $m$ 个未知函数
$y_1,\ldots,y_m$，其中 $y_i$ 的最高阶导数为 $n_i$，并且最高阶导数已经
显式解出. 对每个 $y_i$ 引入状态变量
\[
x_{i,1}=y_i,\quad x_{i,2}=y_i',\quad\ldots,\quad
x_{i,n_i}=y_i^{(n_i-1)}.
\]
每个 $y_i$ 的导数关系给出一部分一阶方程，而原来的最高阶方程给出
\[
x_{i,n_i}'=
f_i\left(t,y_1,y_1',\ldots,x_{1,n_1},\ldots,x_{m,n_m}\right).
\]
把所有 $x_{i,j}$ 按任意固定顺序排成状态向量，便得到一阶微分方程组.
反过来，任意一阶方程组都是这种变换的特殊情形. 因此研究一阶方程组并不失去
一般性.

初始条件也同步变换为
\[
x_i(t_0)=y^{(i-1)}(t_0),\qquad i=1,\ldots,n.
\]
对于含有多个未知函数的高阶方程组，应对每个未知函数逐阶引入状态变量，
直到所有方程只含一阶导数.

\begin{eg}
将
\[
y''-2y'+5y=\sin t,\qquad y(0)=1,\quad y'(0)=0
\]
化为一阶微分方程组.
\begin{pf2}
令 $x_1=y,\ x_2=y'$，则
\[
\left\{
\begin{aligned}
x_1'&=x_2,\\
x_2'&=-5x_1+2x_2+\sin t.
\end{aligned}
\right.
\]
矩阵形式为
\[
\frac{d}{dt}
\begin{pmatrix}x_1\\x_2\end{pmatrix}
=
\begin{pmatrix}0&1\\-5&2\end{pmatrix}
\begin{pmatrix}x_1\\x_2\end{pmatrix}
+\begin{pmatrix}0\\\sin t\end{pmatrix},
\qquad
\begin{pmatrix}x_1(0)\\x_2(0)\end{pmatrix}
=
\begin{pmatrix}1\\0\end{pmatrix}.
\]
\end{pf2}
\end{eg}

\subsection{函数向量、函数矩阵及其范数}

若 $x_i(t)$ 是区间 $I$ 上的函数，则
\[
\bm{x}(t)=(x_1(t),\ldots,x_n(t))^\mathrm{T}
\]
称为\textbf{函数列向量}. 若 $a_{ij}(t)$ 是 $I$ 上的函数，则
\[
A(t)=(a_{ij}(t))_{n\times n}
\]
称为\textbf{函数矩阵}. 函数向量和函数矩阵的加法、数乘、矩阵乘法与常向量、
常矩阵的运算规则相同；其连续性、导数与积分均按分量（元素）定义，即对每个
分量或每个元素分别作相应运算：
\[
\bm{x}'(t)=\bigl(x_1'(t),\ldots,x_n'(t)\bigr)^\mathrm{T},\qquad
A'(t)=\bigl(a_{ij}'(t)\bigr)_{n\times n},
\]
\[
\int_a^b\bm{x}(t)\,dt
=\left(\int_a^b x_1(t)\,dt,\ldots,\int_a^b x_n(t)\,dt\right)^\mathrm{T},
\qquad
\int_a^b A(t)\,dt=\left(\int_a^b a_{ij}(t)\,dt\right)_{n\times n},
\]
只要右端各分量（元素）的导数或积分存在. 当相应导数存在时，乘积的求导仍满足
乘法法则：
\begin{equation}\label{equ:4.6}
\frac{d}{dt}\bigl(A(t)\bm{x}(t)\bigr)
=A'(t)\bm{x}(t)+A(t)\bm{x}'(t).
\end{equation}

对
$\bm{x}=(x_1,\ldots,x_n)^\mathrm{T}$ 和 $A=(a_{ij})_{n\times n}$，定义
\begin{equation}\label{equ:4.7}
\|\bm{x}\|=\sum_{i=1}^n|x_i|,
\qquad
\|A\|=\sum_{i,j=1}^n|a_{ij}|.
\end{equation}

\begin{prop}\label{prop:4.1.2}
上述范数满足
\begin{description}
\item[(i)] $\|\bm{x}\|\ge0,\ \|A\|\ge0$，且范数为零当且仅当向量或矩阵为零；
\item[(ii)] $\|c\bm{x}\|=|c|\|\bm{x}\|,\ \|cA\|=|c|\|A\|$；
\item[(iii)] $\|\bm{x}+\bm{y}\|\le\|\bm{x}\|+\|\bm{y}\|,\
\|A+B\|\le\|A\|+\|B\|$；
\item[(iv)] $\|A\bm{x}\|\le\|A\|\|\bm{x}\|,\
\|AB\|\le\|A\|\|B\|$；
\item[(v)]
$\left\|\int_a^b\bm{x}(s)\,ds\right\|\le\int_a^b\|\bm{x}(s)\|\,ds$.
\end{description}
\end{prop}

\begin{pf1}
除乘积不等式外，其余性质均可由绝对值不等式和有限求和直接得到. 对
$A=(a_{ij})$、$\bm{x}=(x_1,\ldots,x_n)^\mathrm{T}$，有
\[
\|A\bm{x}\|
=\sum_{i=1}^n\left|\sum_{j=1}^n a_{ij}x_j\right|
\le\sum_{i,j}|a_{ij}||x_j|
\le\|A\|\sum_{j=1}^n|x_j|
=\|A\|\|\bm{x}\|.
\]
取 $A,B$ 的乘积后再应用上式即得 $\|AB\|\le\|A\|\|B\|$.
\end{pf1}

若函数向量序列 $\{\bm{x}_k(t)\}$ 的每个分量都在 $I$ 上一致收敛，则称
$\{\bm{x}_k(t)\}$ 在 $I$ 上\textbf{一致收敛}. 设有收敛的数项级数
$\sum M_k$，并使
\[
\|\bm{x}_k(t)\|\le M_k,\qquad t\in I,
\]
则由 Weierstrass 判别法，函数向量级数 $\sum\bm{x}_k(t)$ 在 $I$ 上
一致收敛；当连续函数向量列在闭区间上一致收敛时，还可以交换极限与积分：
\[
\lim_{k\to\infty}\int_{t_0}^{t}\bm{x}_k(s)\,ds
=\int_{t_0}^{t}\lim_{k\to\infty}\bm{x}_k(s)\,ds.
\]

第 \ref{sec:1.2} 节中建立的 Gronwall 不等式（引理 \ref{lem:1.2.1}）在向量情形
下同样适用：只需将其中的函数取为向量函数的范数，便可得到向量版本的估计，
此处不再重复.

\subsection{线性方程组的存在唯一性}

\begin{thm}\label{thm:4.1.1}
设 $A(t)$ 和 $\bm{f}(t)$ 在 $[a,b]$ 上连续，$t_0\in[a,b]$，
$\bm{x}_0\in\R^n$. 则初值问题
\begin{equation}\label{equ:4.8}
\left\{
\begin{aligned}
&\bm{x}'=A(t)\bm{x}+\bm{f}(t),\\
&\bm{x}(t_0)=\bm{x}_0
\end{aligned}
\right.
\end{equation}
在 $(a,b)$ 内存在唯一解.
\end{thm}

\begin{pf1}
连续可微函数 $\bm{x}(t)$ 满足 (\ref{equ:4.8})，当且仅当它满足积分方程
\begin{equation}\label{equ:4.9}
\bm{x}(t)=\bm{x}_0+\int_{t_0}^{t}
\left[A(s)\bm{x}(s)+\bm{f}(s)\right]ds.
\end{equation}
记 $F(t,\bm{x})=A(t)\bm{x}+\bm{f}(t)$. 由于 $A(t),\bm{f}(t)$ 在 $[a,b]$ 上
连续，$F$ 在 $[a,b]\times\R^n$ 上连续，且关于 $\bm{x}$ 满足整体 Lipschitz
条件：
\[
\|F(t,\bm{x})-F(t,\bm{y})\|
\le\|A(t)\|\,\|\bm{x}-\bm{y}\|
\le L\|\bm{x}-\bm{y}\|,
\]
其中 $L=\max\limits_{t\in[a,b]}\|A(t)\|$. 由第 \ref{sec:1.1} 节的存在唯一性
定理（定理 \ref{thm:1.1.1}），初值问题 (\ref{equ:4.8}) 在 $t_0$ 的某邻域内
存在唯一解.

再证解可延拓到整个 $(a,b)$. 对 (\ref{equ:4.8}) 的任一解，由 (\ref{equ:4.9})
得
\[
\|\bm{x}(t)\|
\le\|\bm{x}_0\|+\int_{t_0}^{t}\bigl(L\|\bm{x}(s)\|+K\bigr)ds,
\qquad K=\max_{t\in[a,b]}\|\bm{f}(t)\|.
\]
由 Gronwall 不等式（引理 \ref{lem:1.2.1}），$\bm{x}(t)$ 在 $[a,b]$ 的任意闭
子区间上一致有界，不会在有限时刻发生爆破. 根据解的延拓定理（定理
\ref{thm:1.2.2}），解的存在区间可扩大到 $(a,b)$.
\end{pf1}

若线性方程组为
\[
\bm{x}'=A(t)\bm{x},
\]
则齐次初值问题零初始条件只有零解. 

\section{消元法与首次积分法}\label{sec:4.2}

\subsection{消元法}

对由两个未知函数构成的一阶方程组
\[
x'=f(t,x,y),\qquad y'=g(t,x,y),
\]
可从其中一个方程解出某个未知函数，再求导并代入另一个方程，从而得到一个
只含一个未知函数的高阶方程. 求出该未知函数后，必须利用原方程组回代求另一个
未知函数；若直接对已经求出的关系式求积分，可能引入不属于原方程组的额外常数.

上述消元的关键是不断对关系式求导，再作线性组合. 对于线性方程组，这两类
运算可以用一个记号统一刻画. 记 $D=\frac{d}{dt}$，则 $D^k$ 表示对函数求
$k$ 阶导数；形如
\[
L(D)=D^n+a_1D^{n-1}+\cdots+a_n
\]
（$a_i$ 为常数）的表达式称为\textbf{微分算子多项式}，它作用于函数 $x(t)$
的含义是
\[
L(D)x(t)=x^{(n)}(t)+a_1x^{(n-1)}(t)+\cdots+a_nx(t).
\]
引入算子的好处在于：消元中反复出现的“求导、相加、代入”被翻译成算子多项
式之间的代数运算，因而可以像解普通线性方程组那样，用消元法或行列式
（Cramer 法则）系统地进行. 这一方法成立的基础是：常系数微分算子多项式在
相同函数上可交换，即对任意足够光滑的函数 $x$，
\[
L_1(D)L_2(D)x=L_2(D)L_1(D)x.
\]
若
\[
\left\{
\begin{aligned}
L_1(D)x_1+L_2(D)x_2&=g_1(t),\\
L_3(D)x_1+L_4(D)x_2&=g_2(t),
\end{aligned}
\right.
\]
则用 $L_3(D)$ 作用第一式、用 $L_1(D)$ 作用第二式并相减，得
\begin{equation}\label{equ:4.10}
\bigl(L_1L_4-L_3L_2\bigr)x_2
=L_1g_2-L_3g_1.
\end{equation}
该方程为只含 $x_2$ 的高阶线性方程. 求得 $x_2$ 后，可利用原方程组的一式
求 $x_1$.

\begin{thm}\label{thm:4.2.1}
考虑常系数线性方程组
\begin{equation}\label{equ:4.11}
\sum_{j=1}^{m}L_{ij}(D)x_j=g_i(t),
\qquad i=1,\ldots,m,
\end{equation}
其中 $L_{ij}(D)$ 是常系数微分算子多项式. 记
\[
\Delta(D)=\det\bigl(L_{ij}(D)\bigr)_{m\times m}.
\]
若 $\Delta(D)\not\equiv0$，则每个 $x_j$ 都满足一个仅含 $x_j$ 的高阶
常系数线性方程；这些方程可由 $\Delta(D)$ 与右端函数经有限次求导、相加得到.
\end{thm}

\begin{pf1}
将算子矩阵记为 $\mathcal L(D)$. 由于算子多项式具有交换性，可计算
$\mathcal L(D)$ 的伴随矩阵
$\mathcal L^*(D)$. 有
\[
\mathcal L^*(D)\mathcal L(D)
=\Delta(D)I.
\]
用 $\mathcal L^*(D)$ 左乘 (\ref{equ:4.12}) 的向量形式，即得
\[
\Delta(D)x_j=h_j(t),\qquad j=1,\ldots,m,
\]
其中 $h_j(t)$ 是 $g_1,\ldots,g_m$ 的有限线性组合. 这便得到只含单个未知
函数的高阶方程.
\end{pf1}

\begin{rmk}
用算子消元得到若干个一阶或多阶方程后，仍须代回原方程组检验. 单独求得的
高阶方程可能引入不属于原方程组的解，因此算子消元不能只保留最后得到的标量
方程.
\end{rmk}

还需注意，算子作用于两个函数的乘积时仍须使用乘积法则，不能将 $L(D)$ 当作
普通常数直接处理.

\subsection{首次积分}

\begin{deff}[(首次积分)]
考虑方程组
\begin{equation}\label{equ:4.12}
\frac{dx_i}{dt}=f_i(t,x_1,\ldots,x_n),\qquad i=1,\ldots,n.
\end{equation}
若连续可微函数 $\varphi(t,x_1,\ldots,x_n)$ 不是常数，并且沿
(\ref{equ:4.12}) 的任一解均有
\[
\varphi(t,x_1(t),\ldots,x_n(t))=C,
\]
其中 $C$ 只依赖于所取的具体解，则称
\[
\varphi(t,x_1,\ldots,x_n)=C
\]
为方程组 (\ref{equ:4.12}) 的一个\textbf{首次积分}.
\end{deff}

\begin{thm}\label{thm:4.2.2}
连续可微函数 $\varphi$ 沿方程组 (\ref{equ:4.12}) 的解为常数的充要条件是
\begin{equation}\label{equ:4.13}
\frac{\partial\varphi}{\partial t}
+\sum_{i=1}^{n}f_i(t,x_1,\ldots,x_n)
\frac{\partial\varphi}{\partial x_i}=0.
\end{equation}
\end{thm}

\begin{pf1}
沿任一解 $x_i=x_i(t)$，由复合函数求导得
\[
\frac{d}{dt}\varphi(t,x_1(t),\ldots,x_n(t))
=\varphi_t+\sum_{i=1}^n f_i\varphi_{x_i}.
\]
若 (\ref{equ:4.13}) 成立，则全导数为零，故 $\varphi$ 为常数；反之，若
$\varphi$ 沿任一解都是常数，则其全导数在解经过的区域内恒为零，从而
(\ref{equ:4.13}) 成立.
\end{pf1}

\begin{deff}[(独立性)]
设
\[
\varphi_i(t,x_1,\ldots,x_n)=C_i,\qquad i=1,\ldots,n
\]
为 $n$ 个首次积分. 若
\[
\frac{D(\varphi_1,\ldots,\varphi_n)}
{D(x_1,\ldots,x_n)}
=\det\left(\frac{\partial\varphi_i}{\partial x_j}\right)\ne0,
\]
则称这些首次积分在相应区域内\textbf{互相独立}.
\end{deff}

\begin{thm}\label{thm:4.2.3}
若 $n$ 个未知函数的一阶方程组有 $n$ 个互相独立的首次积分，则可由这
$n$ 个关系式确定方程组的通解. 已知一个首次积分，则可利用它消去一个
未知函数，将原方程组化为含有 $n-1$ 个未知函数的方程组.
\end{thm}

\begin{pf1}
设已知一个首次积分
\[
\varphi(t,x_1,\ldots,x_n)=C_1.
\]
若 $\varphi_{x_n}\ne0$，由隐函数定理可在局部解出
\[
x_n=\psi(t,x_1,\ldots,x_{n-1};C_1).
\]
将它代入前 $n-1$ 个方程，便得到变量
$x_1,\ldots,x_{n-1}$ 的 $n-1$ 个方程. 因此一个首次积分可消去一个状态变量.

若有 $n$ 个互相独立的首次积分，则 $n$ 个方程
\[
\varphi_i(t,x_1,\ldots,x_n)=C_i,\qquad i=1,\ldots,n
\]
关于 $x_1,\ldots,x_n$ 的 Jacobi 行列式不为零. 由隐函数定理可唯一解出
\[
x_i=\psi_i(t;C_1,\ldots,C_n),\qquad i=1,\ldots,n.
\]
这 $n$ 个函数的任意性来自 $n$ 个独立常数，因此构成原方程组的通解.
\end{pf1}

\begin{eg}
求方程组
\[
x'=y,\qquad y'=x
\]
的通解.
\begin{pf2}
两式相加、相减，得
\[
(x+y)'=x+y,\qquad (x-y)'=-(x-y).
\]
因此
\[
(x+y)e^{-t}=C_1,\qquad (x-y)e^t=C_2
\]
是两个首次积分. 其 Jacobi 行列式不为零，故它们独立. 于是
\[
x=\frac12\left(C_1e^t+C_2e^{-t}\right),
\qquad
y=\frac12\left(C_1e^t-C_2e^{-t}\right).
\]
\end{pf2}
\end{eg}

\section{线性微分方程组的基本理论}\label{sec:4.3}

本节讨论
\begin{equation}\label{equ:4.14}
\bm{x}'=A(t)\bm{x}+\bm{f}(t)
\end{equation}
及其对应的齐次方程组
\begin{equation}\label{equ:4.15}
\bm{x}'=A(t)\bm{x}.
\end{equation}
其中 $A(t),\bm{f}(t)$ 在 $[a,b]$ 上连续.

\subsection{齐次方程组的解空间}

\begin{thm}[(叠加原理)]\label{thm:4.3.1}
若 $\bm{x}_1,\ldots,\bm{x}_m$ 是齐次方程组 (\ref{equ:4.15}) 的解，
则任意线性组合
\[
\sum_{i=1}^m C_i\bm{x}_i
\]
仍是 (\ref{equ:4.15}) 的解.
\end{thm}

\begin{pf1}
由 $A(t)$ 的线性性，
\[
\left(\sum_{i=1}^m C_i\bm{x}_i\right)'
=\sum_{i=1}^m C_iA(t)\bm{x}_i
=A(t)\sum_{i=1}^m C_i\bm{x}_i.
\]
\end{pf1}

函数向量组的\textbf{线性相关}与\textbf{线性无关}，与第 \ref{sec:3.2} 节
中标量函数组的相应定义完全类似，只需将标量函数改为函数向量、数 $0$ 改为
零向量 $\bm{0}$，此处不再重复.

\begin{deff}[(Wronskian 行列式)]
以 $n$ 个函数向量 $\bm{x}_1,\ldots,\bm{x}_n$ 为列构造矩阵
\[
\Phi(t)=(\bm{x}_1(t),\ldots,\bm{x}_n(t)),
\]
称
\[
W(t)=\det\Phi(t)
\]
为这些函数向量的 \textbf{Wronskian 行列式}.
\end{deff}

\begin{thm}\label{thm:4.3.2}
齐次方程组 (\ref{equ:4.15}) 的 $n$ 个解
$\bm{x}_1,\ldots,\bm{x}_n$ 在 $[a,b]$ 上线性相关，当且仅当
\[
W(t)\equiv0.
\]
它们线性无关，当且仅当存在 $t_0\in[a,b]$，使 $W(t_0)\ne0$.
\end{thm}

\begin{pf1}
若解组线性相关，则其列向量在任何时刻线性相关，故 $W(t)\equiv0$.

反之，若 $W(t_0)=0$，则存在不全为零的 $C_1,\ldots,C_n$，使
\[
\sum_{i=1}^n C_i\bm{x}_i(t_0)=\bm{0}.
\]
线性组合 $\bm{x}(t)=\sum C_i\bm{x}_i(t)$ 是齐次方程组的解，并满足
$\bm{x}(t_0)=\bm{0}$. 由齐次初值问题解的唯一性，$\bm{x}(t)\equiv\bm{0}$，
故解组线性相关.
\end{pf1}

\begin{thm}[(Liouville 公式)]\label{thm:4.3.3}
齐次方程组的任意 $n$ 个解构成的 Wronskian 行列式满足
\begin{equation}\label{equ:4.16}
W(t)=W(t_0)
\exp\left(\int_{t_0}^{t}\trace A(s)\,ds\right).
\end{equation}
因此 $W(t)$ 在定义区间上要么恒为零，要么恒不为零.
\end{thm}

\begin{pf1}
由行列式求导公式及 $\Phi'=A\Phi$，可得
\[
W'(t)=\trace A(t)\,W(t).
\]
解该一阶线性方程即得 (\ref{equ:4.16}).
\end{pf1}

\begin{rmk}
当高阶线性方程
\[
y^{(n)}+a_1(t)y^{(n-1)}+\cdots+a_n(t)y=f(t)
\]
按命题 \ref{prop:4.1.1} 化为一阶组时，其基解矩阵的行列式恰为第
\ref{sec:3.2} 节中的 Wronskian 行列式；又此时 $\trace A(t)=-a_1(t)$，
故 (\ref{equ:4.16}) 退化为定理 \ref{thm:3.2.9} 的
Liouville 公式. 因而本节的结论是第 \ref{sec:3.2} 节相应结论在矩阵形式
下的推广，而非另起炉灶的重复.
\end{rmk}

\begin{thm}\label{thm:4.3.4}
齐次方程组 (\ref{equ:4.15}) 一定存在 $n$ 个线性无关解，故其解空间是
$n$ 维线性空间.
\end{thm}

\begin{pf1}
任取 $t_0\in[a,b]$，由存在唯一性定理，分别求满足
\[
\bm{x}_i(t_0)=\bm{e}_i,\qquad i=1,\ldots,n
\]
的解 $\bm{x}_i(t)$，其中 $\bm{e}_i$ 为标准单位列向量. 此时
$\Phi(t_0)=I$，故 $W(t_0)=1\ne0$. 由定理 \ref{thm:4.3.2}，这些解
线性无关.
\end{pf1}

\begin{deff}[(基解矩阵)]
齐次方程组 (\ref{equ:4.15}) 的 $n$ 个线性无关解称为\textbf{基本解组}.
以基本解组为列构成的矩阵称为\textbf{基本解矩阵}或\textbf{基解矩阵}.
\end{deff}

\begin{thm}\label{thm:4.3.5}
设 $\Phi(t)$ 是齐次方程组 (\ref{equ:4.15}) 的基解矩阵，则
(\ref{equ:4.15}) 的通解为
\begin{equation}\label{equ:4.17}
\bm{x}(t)=\Phi(t)\bm{C},
\qquad \bm{C}\in\R^n.
\end{equation}
\end{thm}

\begin{pf1}
由叠加原理，$\Phi(t)\bm{C}$ 是解. 反过来，设 $\bm{x}(t)$ 是任一解，
取 $t_0\in[a,b]$，令
\[
\bm{C}=\Phi^{-1}(t_0)\bm{x}(t_0).
\]
则解 $\Phi(t)\bm{C}$ 与 $\bm{x}(t)$ 满足同一初始条件，由唯一性得
$\bm{x}(t)=\Phi(t)\bm{C}$.
\end{pf1}

\begin{thm}\label{thm:4.3.6}
设 $\Phi(t)$ 和 $\Psi(t)$ 都是齐次方程组 (\ref{equ:4.15}) 的基解矩阵，
则存在非奇异常数矩阵 $C$，使
\[
\Psi(t)=\Phi(t)C.
\]
特别地，
\[
\Psi(t)=\Phi(t)\Phi^{-1}(t_0)
\]
是满足 $\Psi(t_0)=I$ 的基解矩阵.
\end{thm}

\begin{pf1}
由于 $\Phi(t)$ 可逆，令
\[
X(t)=\Phi^{-1}(t)\Psi(t).
\]
利用 $\Phi'=A\Phi$、$\Psi'=A\Psi$，得
\[
X'(t)
=-\Phi^{-1}\Phi'\Phi^{-1}\Psi+\Phi^{-1}\Psi'
=-\Phi^{-1}A\Phi\Phi^{-1}\Psi
+\Phi^{-1}A\Psi=0.
\]
故 $X(t)$ 为常矩阵 $C$. 由于两个基解矩阵都可逆，$C$ 也非奇异.
\end{pf1}

\subsection{非齐次线性方程组}

\begin{thm}\label{thm:4.3.7}
设 $\Phi(t)$ 是齐次方程组 (\ref{equ:4.15}) 的基解矩阵，
$\bm{x}_p(t)$ 是非齐次方程组 (\ref{equ:4.14}) 的一个特解，则
(\ref{equ:4.14}) 的通解为
\begin{equation}\label{equ:4.18}
\bm{x}(t)=\Phi(t)\bm{C}+\bm{x}_p(t).
\end{equation}
\end{thm}

\begin{pf1}
任意两个非齐次方程组的解之差满足对应的齐次方程组，因此若
$\bm{x}(t)$ 是任一解，则
$\bm{x}(t)-\bm{x}_p(t)$ 是齐次解，故存在 $\bm{C}$ 使
\[
\bm{x}(t)-\bm{x}_p(t)=\Phi(t)\bm{C}.
\]
反之，(\ref{equ:4.18}) 显然是非齐次方程组的解.
\end{pf1}

\begin{thm}[(常数变易公式)]\label{thm:4.3.8}
设 $\Phi(t)$ 是齐次方程组 (\ref{equ:4.15}) 的基解矩阵，则
(\ref{equ:4.14}) 满足 $\bm{x}(t_0)=\bm{x}_0$ 的解为
\begin{equation}\label{equ:4.19}
\bm{x}(t)
=\Phi(t)\Phi^{-1}(t_0)\bm{x}_0
+\Phi(t)\int_{t_0}^{t}\Phi^{-1}(s)\bm{f}(s)\,ds.
\end{equation}
\end{thm}

\begin{pf1}
令
\[
\bm{x}(t)=\Phi(t)\bm{C}(t).
\]
代入 (\ref{equ:4.14})，并利用 $\Phi'=A\Phi$，得
\[
\Phi(t)\bm{C}'(t)=\bm{f}(t),
\]
故
\[
\bm{C}(t)=\bm{C}(t_0)+\int_{t_0}^{t}\Phi^{-1}(s)\bm{f}(s)\,ds.
\]
由初始条件 $\bm{x}_0=\Phi(t_0)\bm{C}(t_0)$ 即可得到 (\ref{equ:4.19}).
\end{pf1}

\section{常系数齐次线性方程组}\label{sec:4.4}

本节考虑
\begin{equation}\label{equ:4.20}
\bm{x}'=A\bm{x},
\end{equation}
其中 $A$ 为 $n\times n$ 实常数矩阵.

\subsection{特征向量法}

\begin{thm}\label{thm:4.4.1}
设 $A$ 有 $n$ 个互异实特征值 $\lambda_1,\ldots,\lambda_n$，相应的
特征向量为 $\bm{\alpha}_1,\ldots,\bm{\alpha}_n$，则
(\ref{equ:4.20}) 的通解为
\begin{equation}\label{equ:4.21}
\bm{x}(t)=C_1\bm{\alpha}_1e^{\lambda_1t}
+\cdots+C_n\bm{\alpha}_ne^{\lambda_nt}.
\end{equation}
\end{thm}

\begin{pf1}
对每个 $i$，由 $A\bm{\alpha}_i=\lambda_i\bm{\alpha}_i$ 得
\[
\left(\bm{\alpha}_ie^{\lambda_it}\right)'
=\lambda_i\bm{\alpha}_ie^{\lambda_it}
=A\bm{\alpha}_ie^{\lambda_it}.
\]
故各 $\bm{\alpha}_ie^{\lambda_it}$ 均为解. 由于相应特征向量线性无关，
这些解在 $t=0$ 构成基本解组.
\end{pf1}

若实矩阵 $A$ 有共轭复特征值
\[
\lambda=a+bi,\qquad \alpha=\bm{u}+i\bm{v},
\]
则
\[
\bm{z}(t)
=\alpha e^{(a+bi)t}
=\bm{u}(t)+i\bm{v}(t)
\]
是复值解. 与第 \ref{sec:3.3} 节中标量方程的复值解一样（定理 \ref{thm:3.3.1}），
由线性性和实矩阵性质，其共轭、实部与虚部也都是解，故可取实部与虚部得到实解：
\begin{equation}\label{equ:4.22}
\bm{u}(t)
=e^{at}\left(\bm{u}\cos bt-\bm{v}\sin bt\right),
\end{equation}
\begin{equation}\label{equ:4.23}
\bm{v}(t)
=e^{at}\left(\bm{u}\sin bt+\bm{v}\cos bt\right).
\end{equation}

\subsection{重特征值与广义特征向量}

设 $\lambda$ 为 $A$ 的 $k$ 重特征根. 若对应的特征向量个数达到 $k$，
直接取各特征向量对应的解 $\bm{\alpha}_je^{\lambda t}$. 若特征向量个数
不足，则需构造广义特征向量链
\begin{equation}\label{equ:4.24}
(A-\lambda I)\bm{\alpha}_1=0,
\qquad
(A-\lambda I)\bm{\alpha}_j=\bm{\alpha}_{j-1},
\quad j=2,\ldots,k.
\end{equation}

\begin{thm}\label{thm:4.4.2}
设 $\bm{\alpha}_1,\ldots,\bm{\alpha}_k$ 满足
(\ref{equ:4.24})，其中 $\bm{\alpha}_1\ne0$，则
\begin{equation}\label{equ:4.25}
\bm{x}_j(t)
=e^{\lambda t}\left(
\bm{\alpha}_j+t\bm{\alpha}_{j-1}
+\frac{t^2}{2!}\bm{\alpha}_{j-2}
+\cdots
+\frac{t^{j-1}}{(j-1)!}\bm{\alpha}_1
\right),
\quad j=1,\ldots,k
\end{equation}
是 $k$ 个线性无关解.
\end{thm}

\begin{pf1}
记 $N=A-\lambda I$. 对 $\bm{x}_j$ 求导并减去 $A\bm{x}_j$，可得
\[
\bm{x}_j'-A\bm{x}_j
=e^{\lambda t}\left[
\bm{\alpha}_{j-1}-N\bm{\alpha}_j
+t(\bm{\alpha}_{j-2}-N\bm{\alpha}_{j-1})
+\cdots
\right].
\]
由等式链 (\ref{equ:4.24})，每一项均为零，所以
$\bm{x}_j'=A\bm{x}_j$.

下面证明线性无关性. 设
\[
C_1\bm{x}_1(t)+\cdots+C_k\bm{x}_k(t)\equiv0.
\]
在 $t=0$ 得
\[
C_1\bm{\alpha}_1+\cdots+C_k\bm{\alpha}_k=0.
\]
用 $N^{k-1}=(A-\lambda I)^{k-1}$ 作用，并利用
$N^{k-1}\bm{\alpha}_j=0\ (j<k)$、$N^{k-1}\bm{\alpha}_k=\bm{\alpha}_1$，
得到 $C_k\bm{\alpha}_1=0$. 因 $\bm{\alpha}_1\ne0$，故 $C_k=0$.
再依次用 $N^{k-2},\ldots,N$ 作用于 $t=0$ 的关系式，可依次得到
$C_{k-1}=0,\ldots,C_1=0$. 因此这些解线性无关.
\end{pf1}

\begin{eg}
求
\[
\bm{x}'=
\begin{pmatrix}
3&1\\
-1&1
\end{pmatrix}\bm{x}
\]
的通解.
\begin{pf2}
特征方程为
\[
\det(A-\lambda I)=(\lambda-2)^2=0.
\]
取特征向量
\[
\bm{\alpha}_1=
\begin{pmatrix}1\\-1\end{pmatrix}.
\]
解
\[
(A-2I)\bm{\alpha}_2=\bm{\alpha}_1,
\]
可取
\[
\bm{\alpha}_2=
\begin{pmatrix}-1\\0\end{pmatrix}.
\]
因此
\[
\bm{x}_1(t)=
\begin{pmatrix}1\\-1\end{pmatrix}e^{2t},
\qquad
\bm{x}_2(t)=
\left[
\begin{pmatrix}-1\\0\end{pmatrix}
+t\begin{pmatrix}1\\-1\end{pmatrix}
\right]e^{2t}.
\]
通解为
\[
\bm{x}=C_1\bm{x}_1+C_2\bm{x}_2.
\]
\end{pf2}
\end{eg}

\subsection{矩阵指数}

\begin{deff}[(矩阵指数)]
对 $n\times n$ 常数矩阵 $A$，定义
\begin{equation}\label{equ:4.26}
e^{At}
=\sum_{k=0}^{\infty}\frac{A^kt^k}{k!}.
\end{equation}
\end{deff}

\begin{thm}\label{thm:4.4.3}
矩阵指数满足
\begin{description}
\item[(i)] $e^{A0}=I$；
\item[(ii)] $e^{A(t+s)}=e^{At}e^{As}$，从而 $(e^{At})^{-1}=e^{-At}$；
\item[(iii)] $\dfrac{d}{dt}e^{At}=Ae^{At}$；
\item[(iv)] 若 $AB=BA$，则 $e^{A+B}=e^Ae^B$；
\item[(v)] 对任意可逆矩阵 $S$，有
$e^{S^{-1}AS}=S^{-1}e^{AS}S$.
\end{description}
\end{thm}

\begin{thm}\label{thm:4.4.4}
$\Phi(t)=e^{At}$ 是方程组 (\ref{equ:4.20}) 的基解矩阵，且
$\Phi(0)=I$. 因此 (\ref{equ:4.20}) 的通解为
\[
\bm{x}(t)=e^{At}\bm{C}.
\]
\end{thm}

\begin{pf1}
由定理 \ref{thm:4.4.3}，
\[
\frac{d}{dt}e^{At}=Ae^{At},\qquad e^{A0}=I.
\]
由 $A0=0A$ 及矩阵指数的性质得
$e^{A0}=I$，从而行列式为 1，故 $e^{At}$ 是基解矩阵.
\end{pf1}

\begin{cor}\label{cor:4.4.1}
若 $A$ 只有一个特征值 $\lambda$，则
\begin{equation}\label{equ:4.27}
e^{At}
=e^{\lambda t}\sum_{k=0}^{n-1}
\frac{t^k}{k!}(A-\lambda I)^k.
\end{equation}
\end{cor}

\begin{pf1}
令 $N=A-\lambda I$. 由于 $\lambda$ 是唯一的特征值，$N$ 的全部特征值都是
零，故 $N$ 为幂零矩阵. 于是
\[
e^{At}
=e^{\lambda t}e^{Nt}
=e^{\lambda t}\sum_{k=0}^{n-1}\frac{t^k}{k!}N^k.
\]
\end{pf1}

\begin{thm}[(Hamilton--Cayley 定理)]\label{thm:4.4.5}
设 $A$ 的特征多项式为
\[
p_A(\lambda)=\det(\lambda I-A)
=\lambda^n+c_1\lambda^{n-1}+\cdots+c_n,
\]
则
\[
p_A(A)=A^n+c_1A^{n-1}+\cdots+c_nI=O.
\]
因此任意矩阵幂 $A^k\ (k\ge n)$ 都可由
$I,A,\ldots,A^{n-1}$ 线性表示.
\end{thm}

\begin{cor}\label{cor:4.4.2}
对任意实数 $t$，存在函数 $\alpha_0(t),\ldots,\alpha_{n-1}(t)$，使
\begin{equation}\label{equ:4.28}
e^{At}
=\sum_{k=0}^{n-1}\alpha_k(t)A^k.
\end{equation}
\end{cor}

\begin{pf1}
由 $e^{At}$ 的幂级数定义及 Hamilton--Cayley 定理，任意高次幂都可以逐次
降为不超过 $n-1$ 次的矩阵多项式. 有限项合并后即得 (\ref{equ:4.28}).
\end{pf1}

\begin{rmk}
若 $A$ 的特征值 $\lambda_1,\ldots,\lambda_r$ 互异，则可用满足
\[
H_j(\lambda_k)=\delta_{jk}
\]
的 Lagrange 插值多项式 $H_j$，得到
\[
e^{At}=\sum_{j=1}^{r}e^{\lambda_jt}H_j(A).
\]
这给出了不必显式求 $S$ 与 $S^{-1}$ 的一种矩阵指数计算方法.
\end{rmk}

\begin{rmk}
\begin{description}
\item[(i)] 当 $A$ 可对角化时，可写成
$A=S\diag(\lambda_1,\ldots,\lambda_n)S^{-1}$，其中 $S$ 可逆，于是
\[
e^{At}
=S\diag(e^{\lambda_1t},\ldots,e^{\lambda_nt})S^{-1}.
\]
\item[(ii)] 矩阵指数与特征根法本质相同：特征根法直接构造基解矩阵，矩阵指数
给出统一而简洁的算子表示.
\item[(iii)] 若右端出现常数矩阵与函数矩阵的乘积，必须保持乘法顺序，不能将
矩阵当作普通数交换位置.
\end{description}
\end{rmk}

\section{常系数非齐次线性方程组}\label{sec:4.5}

考虑
\begin{equation}\label{equ:4.29}
\bm{x}'=A\bm{x}+\bm{f}(t),
\end{equation}
其中 $A$ 为常数矩阵，$\bm{f}(t)$ 连续.

\subsection{常数变易法}

由定理 \ref{thm:4.3.8}，若取 $\Phi(t)=e^{At}$，则
(\ref{equ:4.29}) 的通解为
\begin{equation}\label{equ:4.30}
\bm{x}(t)
=e^{At}\bm{C}
+\int_{t_0}^{t}e^{A(t-s)}\bm{f}(s)\,ds.
\end{equation}
满足 $\bm{x}(t_0)=\bm{x}_0$ 的解为
\begin{equation}\label{equ:4.31}
\bm{x}(t)
=e^{A(t-t_0)}\bm{x}_0
+\int_{t_0}^{t}e^{A(t-s)}\bm{f}(s)\,ds.
\end{equation}

\subsection{线性变换法}

若 $A$ 有 $n$ 个线性无关特征向量
$\bm{\alpha}_1,\ldots,\bm{\alpha}_n$，令
\[
T=(\bm{\alpha}_1,\ldots,\bm{\alpha}_n),
\qquad \bm{x}=T\bm{y},
\]
则
\begin{equation}\label{equ:4.32}
\bm{y}'
=\diag(\lambda_1,\ldots,\lambda_n)\bm{y}
+T^{-1}\bm{f}(t).
\end{equation}
方程组因此解耦为 $n$ 个一阶线性方程
\begin{equation}\label{equ:4.33}
y_i'=\lambda_iy_i+g_i(t),\qquad i=1,\ldots,n.
\end{equation}
分别求解后再令 $\bm{x}=T\bm{y}$.

\subsection{待定系数法与叠加原理}

第 \ref{sec:3.4} 节针对高阶方程介绍的待定系数法可自然推广到方程组：当右端
为多项式、指数函数与三角函数的有限乘积时，可设出同类形式的特解并比较系数.
与标量情形一致，若指数上的频率恰为系数矩阵的特征值，则需相应抬高特解中多项
式的次数；差别仅在于此处按广义特征向量链的长度衡量重数.

\begin{thm}[(叠加原理)]\label{thm:4.5.1}
若 $\bm{x}_i$ 是非齐次方程组
\[
\bm{x}'=A\bm{x}+\bm{f}_i(t)
\]
的解，则 $\sum\bm{x}_i$ 是
\[
\bm{x}'=A\bm{x}+\sum\bm{f}_i(t)
\]
的解；若 $\bm{x}_1,\bm{x}_2$ 是非齐次方程组 (\ref{equ:4.29}) 的解，
则 $\bm{x}_1-\bm{x}_2$ 是对应齐次方程组的解.
\end{thm}

\begin{thm}\label{thm:4.5.2}
设
\[
\bm{f}(t)=e^{\mu t}\bm{P}_m(t),
\]
其中 $\bm{P}_m(t)$ 是次数不超过 $m$ 的多项式向量. 设 $q$ 为特征值
$\mu$ 对应的最长广义特征向量链长度；若 $\mu$ 不是特征值，则取 $q=0$.
则方程组
\[
\bm{x}'=A\bm{x}+e^{\mu t}\bm{P}_m(t)
\]
存在形如
\begin{equation}\label{equ:4.34}
\bm{x}_p(t)=e^{\mu t}\bm{Q}_{m+q}(t)
\end{equation}
的特解，其中 $\bm{Q}_{m+q}(t)$ 是次数不超过 $m+q$ 的多项式向量.
\end{thm}

\begin{pf1}
将 (\ref{equ:4.34}) 代入方程组，并除以非零因子 $e^{\mu t}$，得到
\[
\bm{Q}'+(\mu I-A)\bm{Q}=\bm{P}_m.
\]
比较各次幂的系数，得到一个含待定向量的线性代数方程组. 当 $\mu$ 不是特征值
时，$\mu I-A$ 可逆，多项式次数不需升高，取 $q=0$ 即可唯一确定系数.
当 $\mu$ 是特征值且存在广义特征向量链时，低次方程之间出现相容性约束；
将多项式次数提高 $q$，即可逐次消除约束并确定全部系数. 因此
(\ref{equ:4.34}) 给出了所求特解的形式.
\end{pf1}

\begin{rmk}
若右端是实正弦、余弦函数，可将其写成复指数函数的实部或虚部，先在复数域中
求特解，再取实部或虚部. 这种处理方法与实值待定系数法结果相同，但能统一
处理不同频率和共振情形.
\end{rmk}

\begin{eg}
求
\[
\bm{x}'=
\begin{pmatrix}
1&2\\
4&3
\end{pmatrix}\bm{x}
+\begin{pmatrix}4\\-2\end{pmatrix}e^t
\]
的一个特解.
\begin{pf2}
矩阵 $A$ 的特征值为 $5,-1$，而 $\mu=1$ 不是特征值，故设
\[
\bm{x}_p=\begin{pmatrix}u\\v\end{pmatrix}e^t.
\]
代入得
\[
u=u+2v+4,\qquad v=4u+3v-2.
\]
解得
\[
u=\frac32,\qquad v=-2.
\]
因此
\[
\bm{x}_p=
\begin{pmatrix}
\frac32\\-2
\end{pmatrix}e^t.
\]
\end{pf2}
\end{eg}

\begin{rmk}
\begin{description}
\item[(i)] 常数变易法适用于连续右端，理论最完整；其代价是通常要计算
矩阵指数或逆矩阵.
\item[(ii)] 线性变换法要求系数矩阵具有足够多的线性无关特征向量；
若矩阵不可对角化，则需先处理广义特征向量链.
\item[(iii)] 待定系数法计算简洁，但只适用于多项式、指数函数、正弦余弦及其
有限乘积型右端.
\item[(iv)] 上述方法讨论的是理论解结构；具体物理、生物或工程建模不在本章范围.
\end{description}
\end{rmk}

